\section*{Exercises about Upper-Triangular Matrices}

\exercise{10}
\begin{xrcs}
  Suppose $T \in \linmap(V)$ and $v_1, \ddd, v_n$ is a basis of $V$. Show that the following are equivalent.
  \begin{enumerate}[label=\textbf{(\alph*)}]
    \item The matrix $\mmatrix(T)$ of $T$ with respect to $v_1, \ddd, v_n$ is lower triangular.
    \item $\myspan{v_k, \dots, v_n}$ is invariant under $T$ for each $k = 1, \ddd, n$.
    \item $T v_k \in \myspan{v_k, \dots, v_n}$ for each $k = 1, \ddd, n$.
  \end{enumerate}

  \begin{xprf}
    Like in \ref{thm: conditions for upper-triangular matrix}, we will prove that (a) $\implies$ (b) $\implies$ (c) $\implies$ (a).

    \StepOne Suppose (a) holds and fix $k \in \{1, \ddd, n \}$. Consider $j \in \{k, \ddd, n\}$ (i.e., $k  \leq j$). The matrix $\mmatrix(T)$ of $T$ looks as follows.

    \begin{minipage}{\linewidth}
      \begin{equation}
        \mathcal{M} (T) =
        \begin{blockarray}{cccccc}
          & v_1     & \cdots & v_j    &   \cdots & v_n     \\
          \begin{block}{c(ccccc)}
            v_1    & A_{1,1} &        &          &        &  \\
            \vdots & \vdots  & \ddots &          &        &  \\
            v_j &    A_{j,1} & \cdots & A_{j,j}  &        &  \\
            \vdots & \vdots  & \ddots & \vdots   & \ddots &  \\
            v_n    & A_{n,1} & \cdots & A_{n,j}  & \cdots & A_{n,n} \\
          \end{block}
        \end{blockarray}.
      \end{equation}
    \end{minipage}

    We have
    \begin{equation}
      Tv_j = A_{j,j} v_{j} + A_{j+1,j} v_{j+1} + \cdots + A_{n,j} v_{n} =  \sum_{\ell=1}^{n} A_{\ell,j} v_\ell.
    \end{equation}

    Therefore, $T v_j \in \myspan{v_j, \ddd, v_n}$. Because
    \begin{equation}
      \myspan{v_j, \ddd, v_n} \subseteq \myspan{v_k, \ddd, v_n},
    \end{equation}

    it's immediate that $Tv_j \in \myspan {v_k, \ddd, v_n}$. Since $v_j$ could be any basis vector in the list $v_k, \ddd, v_n$, $\myspan {v_k, \ddd, v_n}$ is invariant under $T$, completing the proof that (a) implies (b).

    \StepTwo Now, suppose (b) holds. Hence, we assume $\myspan{v_k, \ddd, v_n}$ is invariant under $T$ for each $k = 1, \ddd, n$. In particular,
    \begin{equation}
      T v_j \in \myspan{v_j, \ddd, v_n} \quad \forall j \in \{1, \ddd, n\}.
    \end{equation}

    Thus, (b) implies (c)

    \StepThree Finally, suppose (c) holds. Therefore, $T v_k \in \myspan{v_k, \ldots, v_n}$ for $k = 1, \ddd, n$. This implies that for a fixed $k$, there exist scalars $a_k, \dots, a_n \in \myF$ such that
    \begin{equation}
      \label{eq: lower triangular equation for T v_k}
      T v_k = a_k v_k + \cdots + a_n v_n = 0 v_1 + 0 v_2 + \cdots + 0 v_{k-1} + a_k v_k + \cdots + a_n v_n.
    \end{equation}

    This implies that all the entries above the diagonal of $\mmatrix (T)$ are $0$, because the list $v_1, \ldots, v_n$ is linearly independent and we cannot choose another representation for $T v_k$ in respect to that list. Therefore, the scalars $a_k, \dots, a_n$ coincide with the entries of $\mmatrix(T)$. Thus, $\mmatrix(T)$ is a lower-triangular matrix, completing the proof that (c) implies (a).
  \end{xprf}
\end{xrcs}